\chapter{Inverse-source provenance and concordance}
\label{chap:inverse-reference-appendix}

\section{Source corpus and editorial provenance}
\label{sec:ref-source-provenance}

The canonical volume is a semantic consolidation of six independently
delivered source packages.  They are not numbered editions of one manuscript:
overlap records independent rediscovery or reuse, and disagreement must be
settled mathematically rather than by choosing the newest-looking file.  The
historical paths in \cref{tab:ref-source-packages} were relative to the former
sibling layout under \lean{exponents-and-q-series}; the pinned pre-retirement
revision and repository history preserve them.  The row counts are the
numbers of theorem-like source environments recorded in the machine-readable
concordance.

\begin{table}[htbp]
\centering
\small
\begin{tabularx}{0.98\textwidth}{@{}>{\raggedright\arraybackslash}p{0.34\textwidth}r>{\raggedright\arraybackslash}X@{}}
\toprule
Source package and principal file & Rows & Distinctive contribution and
canonical disposition\\
\midrule
\lean{inverse_q_analog_functions_report/}\newline
\lean{inverse_q_analog_functions.tex}
& 49
& The clearest branchwise narrative, passive-parameter viewpoint, real and
complex sheet atlases, and continuation algorithms.  Its reusable inverse
calculus is absorbed into the universal chapter; its product-specific
geometry is distributed among the finite and infinite product chapters.\\
\addlinespace
\lean{inverse_q_analog_jet_atlas/}\newline
\lean{inverse_q_analog_jet_atlas.tex}
& 38
& Formal triangular reversion, Lagrange--B\"urmann and Lagrange--Good
formulas, Bell-polynomial organization, mixed jets, and coefficient atlases.
These are consolidated, with corrected algebraic hypotheses, in the universal
inversion calculus.\\
\addlinespace
\lean{inverse_q_analogs_extended_report/}\newline
\lean{inverse_q_analogs_extended_report.tex}
& 48
& The broadest collection of exceptional finite regimes, radial inverses,
\(q\)-special functions, certification ideas, and Fabius--Rvachev recovery.
Correct special cases are retained at their natural topical destinations;
overgeneralized claims are corrected or retired.\\
\addlinespace
\lean{inverse_q_analogs_report/}\newline
\lean{inverse_q_analogs_report.tex}
& 59
& Discriminants, remote inverse sheets, real branch geometry, numerical
certification, and the earliest large conjecture register.  Exact results are
deduplicated; symbolic collision claims without independently checkable
certificates remain conditional targets rather than theorems.\\
\addlinespace
\lean{inverse_q_analogs_report-2/}\newline
\lean{inverse_q_analogs_report.tex}
& 29
& An independent synthesis emphasizing reciprocal regimes, roots of unity,
compact coefficient tables, and algorithmic cross-checks.  It supplies
corroborating provenance and occasionally sharper normalizations, not a
presumption of chronological priority.\\
\addlinespace
\lean{q_pochhammer_q_binomial_expansions_report/}\newline
\lean{q_analog_expansions_report.tex}
& 37
& The primary forward-expansion engine: finite and infinite
\(q\)-Pochhammer products, Gaussian and multinomial coefficients,
Bernoulli/polylogarithmic expansions, roots of unity, and \(q\)-special
functions.  Its uniform head--tail remainder arguments are retained wherever
termwise formal expansion alone would be insufficient.\\
\bottomrule
\end{tabularx}
\caption{The six source packages and their canonical roles.}
\label{tab:ref-source-packages}
\end{table}

The editorial disposition is governed by five rules.

\begin{enumerate}
\item A semantic duplicate is stated once, with the strongest correct
  hypotheses and one complete proof.  Several source rows may therefore point
  to one canonical label.
\item A useful source theorem may be strengthened, split into logically
  independent results, or absorbed into a more general lemma.  The
  concordance records the primary destination and explains any one-to-many
  refinement in its disposition note.
\item A false, ill-posed, redundant, or unfalsifiable assertion is not
  preserved merely because it carried a theorem or conjecture heading.
  Interesting statements with genuine missing obligations survive only in
  explicitly labelled \emph{Conjecture} environments.
\item Numerical and symbolic calculations are evidence only to the extent
  certified by reproducible exact arithmetic, interval enclosures, or a proved
  residual bound.  Unsupported large factorizations are certification targets,
  not hidden premises of unconditional theorems.
\item A script, table, captured output, or figure is migrated only when it
  uniquely supports a retained result or reproducible calculation.
  Byte-identical duplicates and generated publication byproducts are not
  separate mathematical sources.  Superseded layouts remain recoverable from
  repository history.
\end{enumerate}

\section{The exhaustive result concordance}
\label{sec:ref-concordance}

The file \lean{theorem_concordance.csv} is the machine-readable ledger for the
source corpus.  It has one header and exactly \(260\) data rows, one for every
environment recognized by the extraction audit.  The package totals are
\[
 49+38+48+59+29+37=260.
\]
By kind, the inventory contains \(131\) theorems, \(28\) propositions,
\(16\) corollaries, \(2\) lemmas, \(32\) conjectures, \(28\) problems or
research problems, \(9\) definitions, \(5\) algorithms, \(4\) examples,
\(4\) principles, and \(1\) computational result.  The extractor found an
intervening \lean{proof} or \lean{proofidea} environment after \(163\) rows
and no such environment after \(97\).  This last field is a syntactic audit
signal, not a verdict on correctness: conjectures and definitions normally
have no proof, while a source proof environment may still contain a gap.

\begin{table}[htbp]
\centering
\small
\begin{tabularx}{0.98\textwidth}{@{}>{\raggedright\arraybackslash}p{0.27\textwidth}>{\raggedright\arraybackslash}X@{}}
\toprule
CSV field group & Meaning\\
\midrule
\lean{source_key}
& Stable key formed from the package name and the source label, or from a
  deterministic kind-and-ordinal key when the source supplied no label.\\
\addlinespace
\lean{source_package}, \lean{source_file}, \lean{source_line}
& Exact source location.  Line numbers belong to the archived source snapshot
  and are provenance coordinates, not canonical page references.\\
\addlinespace
\lean{source_label}, \lean{source_kind}, \lean{source_title},
\lean{source_chapter}, \lean{source_section_path}
& The source manuscript's own classification and surrounding hierarchy.
  These fields preserve what was claimed without endorsing it.\\
\addlinespace
\lean{source_proof_present}
& Whether a \lean{proof} or \lean{proofidea} environment occurs before the
  next recognized result.  It does not assert that the proof is sound,
  complete, or relevant to the exact statement.\\
\addlinespace
\lean{canonical_label}, \lean{canonical_status}
& The primary canonical destination and its strict status.  Repeated labels
  across source rows express semantic deduplication.  A split or partial
  retention is described in \lean{disposition_notes}.\\
\addlinespace
\lean{lean_module}, \lean{lean_declaration}
& An exact formal counterpart, when one has actually been identified and
  validated.  Related supporting lemmas do not justify filling these fields
  for a stronger human theorem.\\
\addlinespace
\lean{disposition_notes}
& The reason for strengthening, merging, splitting, correcting, making
  conditional, retaining as a conjecture, or retiring the source result.\\
\bottomrule
\end{tabularx}
\caption{Interpretation of the concordance columns.}
\label{tab:ref-concordance-fields}
\end{table}

For an assertion, \lean{canonical_status} uses exactly the three-way
vocabulary of \cref{sec:status-vocabulary}:
\emph{Lean-proved}, \emph{human-proved frontier result}, or
\emph{conjecture}.  A definition, notation declaration, or purely descriptive
algorithm is not forced into an assertion status; its retention or retirement
is explained in the disposition field.  Likewise, \emph{merged},
\emph{strengthened}, \emph{corrected}, and \emph{retired} are dispositions,
not proof statuses.

The extractor initially generated the five editorial columns empty, because
text matching cannot decide mathematical equivalence.  In the reviewed
ledger, all \(260\) rows now have a canonical status and disposition note;
\(180\) map to canonical labels, while the remaining \(80\) are explicitly
absorbed, retired, or not separately asserted.  Exact Lean module and
declaration cells remain empty unless an exact formal counterpart is claimed.
The retained reviewed CSV is authoritative.  The revision-backed extractor
\lean{audit/extract_source_results.py} reproduces and verifies its ten
immutable source columns from the commit pinned in
\lean{audit/SOURCE_REVISION}; it refuses to overwrite the five editorial
columns.

\section{Notation and mandatory branch data}
\label{sec:ref-notation}

An inverse symbol denotes a branch only after its data have been fixed.  At
minimum one must record the forward map \(F\), the active variable, all passive
parameters, the target or limiting coordinate, the base point or real domain,
the continuation path or sheet, and every logarithm or fractional-power
determination.  The following table is a compact decoding guide for the
notation used throughout the volume.

\begingroup
\small
\begin{longtable}{@{}>{\raggedright\arraybackslash}p{0.22\textwidth}>{\raggedright\arraybackslash}p{0.29\textwidth}>{\raggedright\arraybackslash}p{0.42\textwidth}@{}}
\caption{Notation and the branch information suppressed by a bare inverse
symbol.}
\label{tab:ref-notation-branches}\\
\toprule
Notation & Meaning and domain & Branch data that must accompany inversion\\
\midrule
\endfirsthead
\toprule
Notation & Meaning and domain & Branch data that must accompany inversion\\
\midrule
\endhead
\bottomrule
\endfoot
\(P_n(a,q)=\QPochhammer{a}{q}{n}\)
& Finite product
  \(\prod_{j=0}^{n-1}(1-aq^j)\), with an empty product equal to one.
& State whether \(a\) or \(q\) is active; fix \(n\) and the other parameter;
  record a base point, real monotonicity interval, or Puiseux sheet.\\
\addlinespace
\(P(a,q)=\QPochhammer{a}{q}{\infty}\)
& Normally convergent product only for \(|q|<1\).  Boundary radial limits and
  exterior reciprocal normalizations are different objects.
& For argument inversion, give the principal interval or the zero-gap index
  and side.  For base inversion, give the interior path or the radial
  normalization and target logarithm.\\
\addlinespace
\(\mathcal A_q(y)\)
& The principal real inverse of
  \(a\mapsto\QPochhammer{a}{q}{\infty}\) for \(0<q<1\), \(y>0\),
  selected by \(a<1\).
& No sheet index is needed on this monotone branch; a nonprincipal real or
  complex inverse must not reuse this symbol without an explicit sheet label.\\
\addlinespace
\(\GaussianBinomial{n}{k}{q}\)
& Gaussian polynomial, including its polynomial continuation where the
  factorial quotient reads \(0/0\).
& Fix \(n,k\); distinguish the positive real branch, a reciprocal
  \(q\leftrightarrow q^{-1}\) branch, and a cyclotomic zero or local sheet.\\
\addlinespace
\(q=\EulerE^{-t}\)
& Logarithmic coordinate near \(q=1\), with \(t=-\PrincipalLogarithm q\).
& Record the logarithm branch and, in a complex asymptotic statement, a
  closed sector contained in \(\Re t>0\).\\
\addlinespace
\(q=\zeta\EulerE^{-t}\)
& Radial or sectorial approach to a primitive root \(\zeta\) of order \(d\).
& Record \(\zeta\), \(d\), the sector or real radial path, and the continuation
  of every logarithm.  This notation never asserts a germ across \(|q|=1\).\\
\addlinespace
\(s=1-y\), \(\eta=((1-a)-y)/(a(1-a))\)
& Regular target coordinates at the unit level and at the small-base endpoint.
& Record the source base point and the nonvanishing denominators.  A different
  normalization changes every printed inverse coefficient.\\
\addlinespace
\(X=-\PrincipalLogarithm y+C\)
& Reciprocal-log scale used when a forward logarithm begins
  \(-A/t+C+\cdots\).
& Fix the target logarithm, the constant \(C\), the sheet with \(t\sim A/X\),
  and a sector on which \(A/X\) is admissible.\\
\addlinespace
\(W_k\)
& The \(k\)-th Lambert-\(W\) determination.
& Record \(k\) and verify the original chosen-log equation.  Solving an
  exponentiated equation alone can admit candidates on the wrong logarithm
  sheet.\\
\addlinespace
\((y-y_0)^{1/m}\)
& Ramified coordinate at a zero of order \(m\).
& Choose the root on a simply connected punctured target domain, or work on
  the ramified cover; all \(m\) determinations coalesce at \(y_0\).\\
\addlinespace
\(\Gamma_q,\psi_q,B_q\)
& The \(q\)-gamma, \(q\)-digamma, and \(q\)-beta normalizations specified in
  the \(q\)-special-function chapter.
& State the real parameter interval or complex continuation, the active
  parameter, the base regime, and any reciprocal continuation for \(q>1\).\\
\addlinespace
\(\varphi_q(t)=\CharacteristicFunction{X_q}(t)\),
\(\CharacteristicFunction{Z_q}(t)\), sinc products
& Characteristic functions of the normalized geometric-uniform law and its
  centered normalization, including the Fabius--Rvachev member.
& Fix the probability normalization, phase convention, real base
  \(0<q<1\), and whether inversion concerns a distribution parameter, a
  Fourier value, or a cumulant observable.\\
\end{longtable}
\endgroup

\section{Source-to-canonical topical crosswalk}
\label{sec:ref-topical-crosswalk}

The row-level concordance is the authoritative provenance mechanism.  For
human navigation, \cref{tab:ref-topical-crosswalk} records where the major
source themes landed.  A source can appear in several canonical chapters, and
a canonical theorem can synthesize several sources.

\begin{table}[htbp]
\centering
\small
\begin{tabularx}{0.98\textwidth}{@{}>{\raggedright\arraybackslash}p{0.24\textwidth}>{\raggedright\arraybackslash}p{0.30\textwidth}>{\raggedright\arraybackslash}X@{}}
\toprule
Source package & Principal source themes & Canonical destinations\\
\midrule
\lean{inverse_q_analog_functions_report}
& Branch orientation, regular inverse germs, sensitivities, real and complex
sheet atlases, continuation.
& \Cref{sec:uq-regular-germs,sec:uq-sensitivities,fp-finite-covering,sec:ip-argument-atlas,sec:certification-frontiers}.\\
\addlinespace
\lean{inverse_q_analog_jet_atlas}
& Triangular coefficient recurrences, Lagrange inversion, Bell polynomials,
mixed jets, multivariate and Good inversion.
& \Cref{sec:uq-formal-inversion,sec:uq-lagrange-burmann,sec:uq-bell-packaging,sec:uq-multivariate-recursion,sec:uq-lagrange-good}.\\
\addlinespace
\lean{inverse_q_analogs_extended_report}
& Exceptional finite bases, \(q=0,\pm1\), radial cyclotomic inversion,
\(q\)-special functions, certification, Fabius--Rvachev recovery.
& \Cref{fp-base-inversion,fp-minus-one-section,sec:ip-real-base,sec:ip-cyclotomic,sec:qs-qgamma,sec:fr-geometric-law,sec:certification-frontiers}.\\
\addlinespace
\lean{inverse_q_analogs_report}
& Discriminants, real gaps and remote sheets, \(q\)-special inverses,
certification algorithms, conjecture register.
& \Cref{fp-finite-covering,sec:ip-argument-atlas,thm:ip-remote-gaps,sec:qs-qdigamma,sec:certification-frontiers,sec:precise-open-frontiers}.\\
\addlinespace
\lean{inverse_q_analogs_report-2}
& Reciprocal regimes, roots of unity, compact inverse tables, independent
algorithmic derivations.
& \Cref{fp-base-inversion,sec:ip-fixed-a-qone,sec:ip-cyclotomic,sec:gaussian-infinity,sec:qs-qexponentials,sec:certification-frontiers}.\\
\addlinespace
\lean{q_pochhammer_q_binomial_expansions_report}
& Exact forward coefficient engines, uniform remainders, Bernoulli and
polylogarithmic expansions, Gaussian/multinomial limits, roots of unity,
\(q\)-gamma and \(q\)-beta.
& \Cref{fp-finite-covering,sec:ip-fixed-a-qone,sec:ip-qx,sec:gaussian-one,sec:qmultinomial,sec:qs-qgamma,sec:qs-qbeta}.\\
\bottomrule
\end{tabularx}
\caption{Topical source-to-canonical crosswalk.}
\label{tab:ref-topical-crosswalk}
\end{table}
\clearpage

\section{Formalization boundary}
\label{sec:ref-formalization-boundary}

The exact declaration-level source map is already given in
\cref{sec:lean-certification-map} and is not duplicated here.  Its current
formal surface has five principal islands:

\begin{enumerate}
\item classical Fabius inverse identities, monotonicity, continuity,
  differentiability, and explicit effective-continuity moduli;
\item the geometric-uniform probability law, its support and self-similarity,
  the geometric sinc characteristic-function identity, certified sinc and CDF
  tails, and an effective forward Fabius modulus;
\item finite \(q\)-Pochhammer dissection and remainder identities, generic
  contracting infinite-product convergence, finite shifts, dissection and
  zeros, and selected Gaussian-polynomial identities including reciprocity
  and exact values at \(q=-1\);
\item Bell/cumulant conversion and a quadratic asymptotic-inversion lemma;
\item real Lambert-\(W\) branch-point geometry: the opposite vertical
  tangents and failure of finite differentiability, together with the signed
  leading square-root equivalences for the principal and lower branches.
\end{enumerate}

The last island comes from the companion Lambert-\(W\) corpus, specifically
\lean{LambertWBranchPointGeometry.lean} and
\lean{LambertWBranchPointAsymptotics.lean}.  It is supporting formalization
for the universal fold discussion, not a seventh inverse-\(q\) precursor;
accordingly it adds no row to the frozen \(260\)-result source ledger.  The
signed leading equivalences do not include a higher Puiseux coefficient or
remainder theorem, still less a formal inversion theorem for a
\(q\)-product.

These declarations support parts of the canonical narrative; they do not
silently formalize stronger neighbors.  In particular, no Lean status is
claimed here for the general analytic and formal inverse theorems beyond the
named Bell and quadratic support.  The fixed-nome convergence, entire-ness,
zero, and simple-derivative subclaims are formalized, but the explicit
uniform-in-nome tails, joint holomorphy, and derivative kernels remain human
proofs, as do remote-gap and log-periodic asymptotics, all-order
\(q\to1\) and cyclotomic expansions, Dedekind-eta completion, Gaussian and
\(q\)-special global inverse atlases, complex certification algorithms, or
the conjectures in
\cref{sec:precise-open-frontiers}.  A theorem in those families remains a
\emph{human-proved frontier result} unless an exact named declaration is
listed for that exact statement.

The Lean table in \cref{sec:lean-certification-map} is explicitly a source
inspection, not a build report.  No Lean process was run while preparing this
appendix, so it contributes no fresh compiler evidence.  Likewise, a
nonempty \lean{lean_module} or \lean{lean_declaration} field in a future
concordance row will be a cross-reference, not by itself proof of a current
aggregate build.  Compiler validation, source inspection, a focused module
build, and a full repository build are distinct evidence levels and must be
reported separately.

\section{Repository and source note}
\label{sec:ref-repository-note}

No citation key occurs in the canonical package at the time of this
consolidation, so this volume does not manufacture a bibliography from the
longer reading lists of its precursors.  Its immediate documentary sources
were the six packages in \cref{tab:ref-source-packages}.  The retained
package-level ledger \lean{PROVENANCE.md}, row-level ledger
\lean{theorem_concordance.csv}, editorial contract \lean{README.md}, and
revision-backed extractor \lean{audit/extract_source_results.py} are the live
canonical records.  Repository history preserves superseded layouts and
retired duplicates.

Classical named tools used in proofs---the inverse-function theorem,
Lagrange--B\"urmann and Lagrange--Good inversion, Rouch\'e's theorem,
Dedekind-eta modularity, and standard special-function identities---are
invoked with their exact hypotheses at the point of use.  If external
citations are introduced in a later edition, their actual citation keys and
bibliographic entries should be added together; an uncited generic reading
list would not improve the present provenance record.
